\Needspace{7\baselineskip}
\pdfbookmark[0]{Anzeigen. Müller: Theorie der Parallelen. 1822.}{anzeigen-mueller-1822}
\begin{center}
\rule{.85\textwidth}{.3pt}\par
Göttingische gelehrte Anzeigen.\quad 28. Oktober 1822.\par
\rule{.85\textwidth}{.3pt}
\end{center}
\hspace*{\parindent}\textit{Theorie der Parallelen, von} \textsc{Carl Reinhard Müller}, \textit{Doktor der Philosophie, außerordentlichem Professor der Mathematik usw.} 40 Seiten im Quartformat. Marburg 1822.

Der Rezensent hat schon vor sechs Jahren in diesen Blättern seine Überzeugung ausgesprochen, dass alle bisherigen Versuche, die Theorie der Parallelen streng zu beweisen oder die Lücke in der \textsc{Euklid}ischen Geometrie zu schließen, uns diesem Ziel nicht nähergebracht haben. Er kann dieses Urteil nur auf alle späteren Versuche ausdehnen, die ihm bekannt geworden sind. Dennoch bleiben manche dieser Versuche, obwohl sie ihr eigentliches Hauptziel verfehlen, wegen des darin gezeigten Scharfsinns für Freunde der Geometrie lesenswert, und. Der Rezensent glaubt in dieser Hinsicht die vorliegende, anlässlich einer Schulprüfung veröffentlichte kleine Schrift besonders hervorheben zu müssen. Den ganzen einfallsreichen Gedankengang des Verfassers hier ausführlich darzulegen, würde für unsere Blätter zu weit führen und wäre auch überflüssig, da die Schrift selbst gelesen zu werden verdient. Sie hat aber wie alle übrigen Versuche eine schwache Stelle, und sie herauszuarbeiten ist der Zweck dieser Anzeige. Wir finden diese schwache Stelle auf S.~15 im Beweis des Lehrsatzes von Artikel~15. Dieser Satz ist der eigentliche Nerv der ganzen Theorie, die fällt, sobald sich der Satz nicht streng beweisen lässt. Wir geben deshalb zunächst den Lehrsatz wieder. Die zugehörige Figur kann jeder leicht selbst zeichnen.

Wenn jeder Winkel an der Grundseite \(ON\) eines gleichschenkligen Dreiecks größer ist als der Winkel an der Spitze \(A\) und man in \(O\) an die Seite \(OA\) einen Winkel von der Größe des Winkels \(A\) anträgt, dessen anderer Schenkel \(OL\) die Gerade \(AN\) im Punkt \(L\) zwischen \(A\) und \(N\) trifft, und wenn man dann
